\documentclass[twocolumn,pra,amsmath,amssymb,superscriptaddress]{revtex4-2}

\usepackage{graphicx}
\usepackage{hyperref}
\usepackage{booktabs}
\usepackage{xcolor}

\begin{document}

\title{The Bargmann--Segal Lattice:\\
$\pi$-Free Discretization of the Harmonic Oscillator on $S^5$}

\author{J.~Loutey}
\affiliation{Independent Researcher, Kent, Washington}

\date{April 10, 2026}

\begin{abstract}
We construct a discrete graph encoding of the three-dimensional
isotropic harmonic oscillator on the holomorphic sector of $S^5$ and
verify by exact rational arithmetic that the resulting graph and its
projection to physical energies contain no transcendental constants.
The Bargmann--Segal transform~\cite{bargmann1961, hall1994} maps the
harmonic oscillator Hamiltonian to $\hbar\omega(\hat{N} + 3/2)$ on the
Fock--Bargmann space, where $\hat{N}$ is the Euler number operator on
holomorphic monomials. Restricting to the unit sphere $S^5 \subset
\mathbb{C}^3$ identifies the eigenspaces with the symmetric SU(3)
irreps $(N, 0)$, of dimension $(N+1)(N+2)/2$, which match the harmonic
oscillator shell degeneracies exactly.  The graph constructed from
these states with dipole transition weights computed by exact
\texttt{fractions.Fraction} arithmetic is bit-exactly $\pi$-free at
$N_{\max} = 5$ (56 nodes, 165 edges, zero irrationals encountered).
The single $\pi$ in the entire ecosystem is the Gaussian normalization
$\pi^{-3}$ of the continuous Bargmann measure, which the discrete
graph never invokes.  The Coulomb problem on $S^3$ and the harmonic
oscillator on $S^5$ are both natural sphere quotients, but they live
on different spheres with different geometric character (Riemannian vs
complex), have natural operators of different order (second vs first),
yield spectra of different growth (quadratic vs linear), and demand
projections of different type (nonlinear vs linear-affine).  This
structural asymmetry refines the exchange constant
taxonomy~\cite{paper18}: calibration $\pi$ is Coulomb-specific. It is
the transcendental cost of mapping a quadratic Riemannian eigenvalue
to an inverse-square physical energy, and it has no analog for the
harmonic oscillator's linear complex-analytic projection.
\end{abstract}

\maketitle

%================================================================
\section{Introduction}
%================================================================

Fock's 1935 stereographic projection~\cite{fock1935} maps the
hydrogen Schr\"odinger equation onto the free Laplacian on the
three-sphere $S^3$.  Paper~0~\cite{paper0} discretized this
construction: nodes labeled by quantum numbers $(n, l, m)$, edges
from angular-momentum selection rules, and a graph Laplacian whose
eigenvalues converge to those of the continuous Laplace--Beltrami
operator.  The kinetic constant $\kappa = -1/16$ then maps the
discrete eigenvalues $n^2 - 1$ to physical Rydberg energies
$E_n = -Z^2 / (2 n^2)$.  Paper~7~\cite{paper7} validated this
$S^3$ reduction with eighteen independent symbolic proofs.
Paper~18~\cite{paper18} identified a transcendental cost in the
projection: $\pi$ enters through the spectral zeta values of $S^3$
as a \emph{calibration} exchange constant.

The harmonic oscillator $V(r) = \frac{1}{2} m \omega^2 r^2$ has a
different dynamical symmetry group (SU(3) instead of SO(4)) and a
qualitatively different spectrum (linear in shell index instead of
inverse-square).  The natural question is whether an analogous
discretization exists for the harmonic oscillator, and if so, what
role $\pi$ plays.

We answer both questions in this paper.  The discretization exists,
via the Bargmann--Segal transform~\cite{bargmann1961, hall1994}
applied to the Hardy space $H^2(S^5)$.  The graph is bit-exactly
$\pi$-free at every finite $N_{\max}$.  The single $\pi$ in the
entire construction appears in the Gaussian normalization $\pi^{-3}$
of the continuous Bargmann measure, which is never used by the
discrete graph.  We further show that the Coulomb / harmonic
oscillator asymmetry has a structural origin in the difference
between Riemannian and complex natural geometries, and that this
asymmetry refines the exchange constant taxonomy of Paper~18:
\emph{calibration} exchange constants are tied to second-order
Riemannian operators with nonlinear projections, while linear
projections from first-order complex-analytic operators introduce no
calibration constants of any kind.  The result has direct
implications for the framing of the fine-structure constant
conjecture~\cite{paper2_alpha} and for the precision of nuclear
shell-model qubit Hamiltonians built on harmonic oscillator
single-particle bases~\cite{paper23}.

The plan of the paper is as follows.  Section~\ref{sec:bs}
summarizes the Bargmann--Segal construction and its identification
with the SU(3) symmetric representation theory.
Section~\ref{sec:graph} constructs the graph and states the
$\pi$-free certificate (Theorem~\ref{thm:pi-free}).
Section~\ref{sec:asymmetry} presents the Coulomb / HO structural
asymmetry as the conceptual centerpiece (Theorem~\ref{thm:asymmetry}).
Section~\ref{sec:rigidity} states the harmonic oscillator rigidity
theorem (Theorem~\ref{thm:rigidity}), the dual of the Fock projection
rigidity theorem of Paper~23~\cite{paper23}.  Section~\ref{sec:imp}
discusses implications for the $\alpha$ conjecture, for nuclear
quantum simulation, and for the framework's universal /
Coulomb-specific partition.

%================================================================
\section{The Bargmann--Segal construction}
\label{sec:bs}
%================================================================

\subsection{The Bargmann transform}

The Bargmann transform~\cite{bargmann1961} is the integral operator
\begin{equation}
  [B f](z) =
  \pi^{-3/4}
  \int_{\mathbb{R}^3}
  \exp\!\Bigl(
    -\tfrac{1}{2}(z^2 + x^2) + \sqrt{2}\, z \cdot x
  \Bigr)\, f(x)\, d^3 x,
  \label{eq:bargmann-kernel}
\end{equation}
mapping $L^2(\mathbb{R}^3)$ unitarily onto the
Fock--Bargmann space $F^2(\mathbb{C}^3)$ of holomorphic functions on
$\mathbb{C}^3$ that are square-integrable against the Gaussian
measure
\begin{equation}
  d\mu(z) = \pi^{-3}\, e^{-|z|^2}\, d^6 z.
  \label{eq:bargmann-measure}
\end{equation}
The harmonic-oscillator eigenstates $|N, l, m\rangle$ map under $B$
to solid harmonic polynomials of total degree $N$ in the complex
coordinates $z_1, z_2, z_3$.  The transform is unitary, so the
eigenvalue equations carry across exactly.

\subsection{The Euler operator}
\label{sec:euler}

The number (Euler) operator on $F^2(\mathbb{C}^3)$ is
\begin{equation}
  \hat{N} = \sum_{i = 1}^{3} z_i\, \frac{\partial}{\partial z_i},
  \label{eq:euler}
\end{equation}
which acts on the monomial $z_1^a z_2^b z_3^c$ with eigenvalue
$a + b + c = N$.  By multinomial counting, the eigenspace
$\hat{N} = N$ has dimension
\begin{equation}
  \dim\bigl\{ z_1^a z_2^b z_3^c \,:\, a + b + c = N \bigr\}
  \;=\;
  \binom{N + 2}{2}
  \;=\;
  \frac{(N+1)(N+2)}{2},
  \label{eq:dim}
\end{equation}
which matches the 3D harmonic oscillator shell degeneracy and the
Weyl dimension formula for the symmetric irreducible representation
$(N, 0)$ of SU(3).

Under the Bargmann transform, the harmonic oscillator Hamiltonian
maps to
\begin{equation}
  H_{\mathrm{HO}}
  \;=\;
  -\tfrac{1}{2 m}\nabla^2 + \tfrac{1}{2} m \omega^2 r^2
  \;\stackrel{B}{\longleftrightarrow}\;
  \hbar \omega \bigl(\hat{N} + \tfrac{3}{2}\bigr),
  \label{eq:H-bargmann}
\end{equation}
so the spectrum is
\begin{equation}
  E_N \;=\; \hbar\omega\bigl(N + \tfrac{3}{2}\bigr)
  \label{eq:HO-spectrum}
\end{equation}
on the eigenspace of dimension $(N+1)(N+2)/2$.  The HO Hamiltonian
is a linear affine function of the Euler operator; the projection
from the discrete eigenvalue to the physical energy involves only
the universal scale $\hbar\omega$ and the additive zero-point
shift $3/2$.

The unitary equivalence of the Schr\"odinger and Bargmann pictures
is rigorous and originates with Bargmann~\cite{bargmann1961} and
Segal~\cite{segal1963}; the modern functional-analytic treatment is
due to Hall~\cite{hall1994}, who extended the construction to
arbitrary compact Lie groups.

\subsection{The $S^5$ Hardy space and SU(3) symmetric representations}

The unit sphere
$S^5 = \{ z \in \mathbb{C}^3 \,:\, |z| = 1 \}$
carries a Cauchy--Riemann (CR) structure inherited from
$\mathbb{C}^3$.  The Hardy space $H^2(S^5)$ consists of boundary
values of holomorphic functions on the unit ball $B^3 \subset
\mathbb{C}^3$ that are square-integrable on $S^5$ with respect to
the surface measure.  Restriction to $S^5$ identifies $H^2(S^5)$
with the closure of the polynomial algebra $\mathbb{C}[z_1, z_2,
z_3]$, with the natural SU(3) action by linear transformations of
the coordinates.

The decomposition by polynomial degree gives an orthogonal direct
sum
\begin{equation}
  H^2(S^5)
  \;=\;
  \bigoplus_{N = 0}^{\infty}
  \mathcal{H}_N,
  \qquad
  \dim \mathcal{H}_N = \frac{(N+1)(N+2)}{2},
  \label{eq:hardy-decomp}
\end{equation}
where $\mathcal{H}_N$ is the space of homogeneous degree-$N$
holomorphic polynomials.  Each $\mathcal{H}_N$ transforms as the
SU(3) irreducible representation $(N, 0)$ — the symmetric tensor
representation of rank $N$ on the fundamental $\mathbf{3}$.

By the Borel--Weil theorem~\cite{borelweil}, this is the same data
as the space of holomorphic sections of the line bundle
$\mathcal{O}(N)$ over $\mathbb{CP}^2$, where $\mathbb{CP}^2 =
\mathrm{SU}(3) / \mathrm{U}(2)$ is the natural quotient.  The
fibration
\begin{equation}
  S^1 \;\longrightarrow\; S^5 \;\longrightarrow\; \mathbb{CP}^2
  \label{eq:s5-fibration}
\end{equation}
is the structural analog of the Hopf fibration $S^1 \to S^3 \to
S^2$ that played a central role in Paper~7~\cite{paper7} and the
$\alpha$-conjecture of Paper~2~\cite{paper2_alpha}.

The CR Laplacian (Kohn Laplacian) $\Box_b$ on the holomorphic
sector of $S^5$, the Dolbeault Laplacian on $\mathcal{O}(N) \to
\mathbb{CP}^2$, and the Euler operator on the Bargmann monomials
all act on the same SU(3) irreducibles $(N, 0)$ and produce
spectra that are linear in the shell index $N$ up to additive and
multiplicative normalizations~\cite{follandstein1974}.  For the
discretization in Section~\ref{sec:graph} we use the Bargmann
description, which is the most elementary: the Euler operator is
diagonal on the monomial basis, requires no differential geometry,
and is amenable to immediate exact arithmetic.

%================================================================
\section{The Bargmann graph}
\label{sec:graph}
%================================================================

\subsection{Graph construction}

\paragraph{Nodes.}
The nodes of the Bargmann graph at cutoff $N_{\max}$ are triples
$(N, l, m_l)$ with $N = 0, 1, 2, \ldots, N_{\max}$; for each $N$,
the orbital angular momentum runs over $l = N, N - 2, \ldots, 1$
or $0$ (same parity as $N$); and for each $l$, the magnetic
quantum number runs over $m_l = -l, \ldots, +l$.  This is the
SU(3) $\to$ SO(3) branching of the symmetric irrep $(N, 0)$.  The
total node count is
\begin{equation}
  |V(N_{\max})|
  \;=\;
  \sum_{N = 0}^{N_{\max}}
  \frac{(N+1)(N+2)}{2}.
  \label{eq:node-count}
\end{equation}
At $N_{\max} = 5$ this gives $1 + 3 + 6 + 10 + 15 + 21 = 56$
nodes.

\paragraph{Edges.}
Edges connect states under SU(3) dipole selection rules:
\begin{equation}
  \Delta N = \pm 1,
  \qquad
  \Delta l = \pm 1,
  \qquad
  \Delta m_l = 0, \pm 1.
\end{equation}
These are the matrix elements of the position operator $z_q$
($q = 0, \pm 1$ in spherical components) between adjacent shells,
which raise $N$ by one and change $l$ by $\pm 1$ by the standard
angular-momentum dipole rule.

\paragraph{Edge weights.}
For each allowed transition the edge weight is the squared matrix
element
\begin{equation}
  w_{(N,l,m_l) \to (N+1,l',m_l')}
  \;=\;
  R^{2}_{n_r,\, l \to l'}
  \;\times\;
  \bigl|
  C^{l' m_l'}_{l\, m_l;\, 1\, q}
  \bigr|^2,
  \label{eq:edge-weight}
\end{equation}
where the radial reduced matrix elements (in HO units $\hbar = m =
\omega = 1$) are
\begin{align}
  \bigl|\langle n_r, l + 1 \,|\, r \,|\, n_r, l \rangle\bigr|^2
  &= n_r + l + \tfrac{3}{2}
  = \tfrac{2 n_r + 2 l + 3}{2},
  \\[2pt]
  \bigl|\langle n_r + 1, l - 1 \,|\, r \,|\, n_r, l \rangle\bigr|^2
  &= n_r + 1,
\end{align}
and the squared rank-1 Wigner $3j$ coefficients factorize into
elementary rationals via the closed-form expressions of
Edmonds~\cite{edmonds} and Varshalovich~\cite{varshalovich}.  Both
factors in Eq.~\eqref{eq:edge-weight} are exact rational numbers,
so the edge weight is exact rational.

\paragraph{Implementation.}
The reference implementation in
\texttt{geovac/nuclear/bargmann\_graph.py} uses Python's
\texttt{fractions.Fraction} type throughout the construction.  No
floating-point arithmetic is used in the build path; the squared
$3j$ tables are evaluated symbolically from the closed-form
expressions, and the radial reduced matrix elements are constructed
directly from the integer formulas above.  Conversion to floating
point occurs only when the dense Hamiltonian or adjacency is
required for numerical eigensolvers.

\subsection{The $\pi$-free certificate}
\label{sec:pi-free}

\begin{theorem}[Bargmann graph rationality]
\label{thm:pi-free}
Let $G(N_{\max})$ be the Bargmann graph constructed in
Section~\ref{sec:graph} with diagonal entries
\begin{equation}
  D_v
  \;=\;
  \mathrm{Fraction}\!\left( 2 N_v + 3,\, 2 \right),
\end{equation}
in units of $\hbar\omega$, and adjacency weights given by
Eq.~\eqref{eq:edge-weight} computed in exact rational arithmetic.
Then every diagonal entry and every nonzero adjacency weight is a
rational number, and no transcendental constant appears in the
construction at any finite $N_{\max}$.
\end{theorem}

\begin{proof}[Proof.]
The diagonal entries are $\mathrm{Fraction}(2 N + 3, 2)$ by
construction (Section~\ref{sec:euler}, Eq.~\eqref{eq:HO-spectrum}).
The adjacency weights are products of two rational factors:
the squared radial reduced matrix elements (which are rational
linear functions of the integer quantum numbers $n_r$ and $l$) and
the squared rank-1 Wigner $3j$ coefficients (which are rational
quotients of polynomials in the integer quantum numbers $l$, $l'$,
$m_l$, $m_l'$ by the Edmonds~\cite{edmonds} formulas).  The product
of two rationals is rational, and the construction terminates after
finitely many additions and multiplications of rationals at any
finite $N_{\max}$.  $\hfill\square$
\end{proof}

\paragraph{Numerical verification.}
An independent run of the function
\texttt{verify\_pi\_free(N\_max=5)} in the reference
implementation returns the certificate:
\begin{quote}\small
\texttt{N\_max: 5}\\
\texttt{n\_nodes: 56}\\
\texttt{n\_edges: 165}\\
\texttt{all\_diagonal\_rational: True}\\
\texttt{all\_adjacency\_rational: True}\\
\texttt{irrationals\_encountered: []}\\
\texttt{pi\_free: True}
\end{quote}
The 17 supporting unit tests are listed in
\texttt{tests/test\_bargmann\_graph.py} and cover node enumeration,
shell degeneracies, parity of $l$ within each shell, the diagonal
HO spectrum, the dipole selection rules, Hermiticity of the
adjacency, exact rationality of every entry, and graph Laplacian
positivity.  All 17 pass.

\paragraph{Where the only $\pi$ lives.}
The single $\pi$ in the entire ecosystem appears in the Gaussian
normalization $\pi^{-3}$ of the continuous Bargmann measure
Eq.~\eqref{eq:bargmann-measure}.  In Paper~18's exchange constant
taxonomy~\cite{paper18}, this is an \emph{embedding} exchange
constant: it quantifies the cost of representing the Bargmann
Hilbert space as a measure space in $6$ real dimensions, and it is
fixed by the requirement that the Gaussian have unit total integral.
The discrete graph never invokes this normalization.  No analogous
$\pi$ enters the spectrum, the degeneracies, the selection rules,
the Wigner $3j$ coefficients, or the radial matrix elements.

\subsection{Hamiltonian and graph Laplacian roles}
\label{sec:roles}

The HO Hamiltonian on the Bargmann graph is a diagonal matrix:
\begin{equation}
  H_{\mathrm{HO}}
  \;=\;
  \hbar\omega \,
  \mathrm{diag}\!\left( N_v + \tfrac{3}{2} \right)_{v \in V}.
  \label{eq:H-graph}
\end{equation}
The eigenvalues are $\hbar\omega(3/2, 5/2, 7/2, 9/2, \ldots)$
with multiplicities $1, 3, 6, 10, \ldots = (N+1)(N+2)/2$.  No
off-diagonal terms are needed to recover the spectrum: it is
simply the diagonal.

The adjacency matrix $A$ encodes dipole transition amplitudes
between HO states (electromagnetic absorption and emission lines).
It plays no role in the construction of the spectrum; its role is
to describe the connectivity structure that determines selection
rules for transition processes.

This is qualitatively different from the Coulomb / $S^3$ case of
Paper~0~\cite{paper0} and Paper~7~\cite{paper7}, where the graph
Laplacian $D - A$ on the $(n, l, m)$ lattice has eigenvalue $n^2 -
1$ on the $n$-th shell and the universal kinetic constant $\kappa
= -1/16$ converts these to Rydberg energies.  There the
off-diagonal structure of $A$ is essential to the spectrum.  In the
Bargmann graph, by contrast, the off-diagonal structure is
spectroscopically informative but spectrally inert.  This is the
first key structural difference between the Coulomb and the
harmonic oscillator discretizations, and it motivates the
asymmetry analysis of Section~\ref{sec:asymmetry}.

%================================================================
\section{The Coulomb / HO asymmetry}
\label{sec:asymmetry}
%================================================================

The Bargmann graph and the Coulomb $S^3$ graph are both natural
discretizations of central-potential quantum mechanics, but they
differ in every structural feature except the use of angular
momentum as a labeling system.  Table~\ref{tab:asymmetry}
summarizes the comparison.

\begin{table*}[t]
\centering
\caption{Structural comparison of the Coulomb and harmonic oscillator
discretizations.  Each row is a feature; each column is a system.
The asymmetry has a causal structure: the geometry choice (row 2)
determines the operator order (row 3), which determines the
spectrum growth (row 4), which determines the projection type
(row 6), which determines whether transcendental constants are
needed (row 9).}
\label{tab:asymmetry}
\begin{ruledtabular}
\begin{tabular}{lll}
Property & Coulomb ($S^3$, Riemannian) & HO ($S^5$ Hardy, complex) \\
\hline
Symmetry group & $\mathrm{SO}(4)$ & $\mathrm{SU}(3)$ \\
Manifold & $S^3$ & $S^5$ (holomorphic sector) \\
Natural operator & Laplace--Beltrami (2nd order) & Euler / Dolbeault (1st order) \\
Eigenvalue $\lambda_N$ & $n^2 - 1$ (quadratic in $n$) & $N$ (linear in $N$) \\
Shell degeneracy & $n^2$ & $(N+1)(N+2)/2$ \\
Projection $E(\lambda)$ & $E_n = -Z^2 / [2(\lambda + 1)]$ (nonlinear) & $E_N = \hbar\omega(\lambda + 3/2)$ (linear-affine) \\
Projection constant & $p_0 = Z / n$ (state-dependent) & $\hbar\omega$ (universal) \\
Graph Laplacian role & computes spectrum via $D - A$ & encodes transitions only \\
Transcendentals in projection & $\pi$ (calibration exchange constant) & none \\
$\pi$-free graph & yes & yes \\
$\pi$ in full pipeline & enters at projection & never enters \\
\end{tabular}
\end{ruledtabular}
\end{table*}

\subsection{The causal chain}
\label{sec:causal}

The asymmetry is not a list of independent observations.  It has a
single causal structure in which each step is forced by the
previous one.

\begin{enumerate}
\item \emph{Geometry choice.}  $\mathrm{SO}(4)$ symmetry of the
  Coulomb problem (the Runge--Lenz vector closes with the angular
  momenta into $\mathfrak{so}(4)$) makes $S^3$ the natural
  manifold, by Fock's reduction~\cite{fock1935}.  $\mathrm{SU}(3)$
  symmetry of the harmonic oscillator (the ladder operators
  $a_i = (x_i + i p_i / m\omega) \sqrt{m\omega / 2\hbar}$
  transform as the fundamental $\mathbf{3}$ and generate
  $\mathfrak{su}(3)$ under commutation) makes $\mathbb{C}^3$ and
  its quotients $\mathbb{CP}^2$ and $S^5$ natural,
  by~\cite{jauchhill1940}.

\item \emph{Operator order.}  A real Riemannian metric on $S^3$
  makes the scalar Laplace--Beltrami operator (second-order in
  derivatives) the natural object.  A complex Cauchy--Riemann or
  K\"ahler structure on $S^5$ or $\mathbb{CP}^2$ makes the
  Dolbeault, Kohn, or Euler operators (first-order in $\hat{N}$
  on holomorphic sections) natural.  This is not a free choice;
  it follows from the requirement that the operator be invariant
  under the dynamical symmetry group acting holomorphically on
  the manifold's structure.

\item \emph{Spectrum growth.}  A second-order self-adjoint
  operator on a compact Riemannian manifold has eigenvalues that
  grow polynomially with shell index, typically as $n^2$ for
  round spheres (Weyl asymptotics).  A first-order equivariant
  operator on the symmetric irreducible representations of a
  compact Lie group has eigenvalues linear in the highest weight.
  These growths are different by definition.

\item \emph{Projection type.}  Mapping a quadratic mathematical
  eigenvalue $\lambda_n = n^2 - 1$ to an inverse-square physical
  energy $E_n = -Z^2 / (2 n^2)$ is a \emph{nonlinear} operation:
  $E = -Z^2 / [2(\lambda + 1)]$.  Mapping a linear mathematical
  eigenvalue $\lambda_N = N$ to a linear physical energy $E_N =
  \hbar\omega(N + 3/2)$ is a \emph{linear-affine} operation.

\item \emph{Spectral zeta and boundary corrections.}  A nonlinear
  projection on a compact manifold requires regularization at the
  boundary of the spectrum, captured by spectral zeta values
  $\zeta_M(s) = \sum_n \lambda_n^{-s}$.  For round $S^d$ these
  involve $\pi$ through Bernoulli numbers and Riemann zeta values
  at integer arguments.  A linear-affine projection has no such
  regularization: every term in the spectrum is finite, the sum
  is unconditionally convergent (or trivially conditionally
  convergent with the additive constant absorbed), and no
  transcendental seed appears.

\item \emph{Calibration exchange constant.}  The transcendentals
  in the spectral zeta are inherited by the physical spectrum as
  what Paper~18~\cite{paper18} terms \emph{calibration} exchange
  constants.  They are the cost of going from the dimensionless
  graph to a physical energy via the nonlinear projection.  For
  the harmonic oscillator this cost is zero; for the Coulomb
  problem it is $\pi$.
\end{enumerate}

We summarize this causal chain as a structural observation:

\begin{theorem}[Structural origin of calibration $\pi$]
\label{thm:asymmetry}
A quantum system whose spectrum is encoded by a self-adjoint
graph Laplacian on a compact Riemannian manifold via a nonlinear
projection requires a calibration exchange constant arising from
the spectral zeta values of the manifold.  A quantum system
whose spectrum is encoded by a first-order equivariant operator
on a compact complex manifold via a linear-affine projection
requires no such calibration constant.  The Coulomb potential and
the three-dimensional isotropic harmonic oscillator instantiate
the two cases respectively.
\end{theorem}

This is a structural description rather than a stand-alone
theorem in the strict mathematical sense.  Each link in the
causal chain is established by independent classical results
(Weyl asymptotics; Borel--Weil; Folland--Stein; Paper~18); the
contribution of the present paper is to assemble them into a
single statement that the role of $\pi$ in the Coulomb projection
is structurally tied to second-order Riemannian geometry, and
that there is no analogous role for $\pi$ in the harmonic
oscillator's complex-analytic projection.

\subsection{Refinement of the exchange constant taxonomy}
\label{sec:taxonomy}

Paper~18~\cite{paper18} classified exchange constants into four
types: \emph{intrinsic} (internal to the discrete graph),
\emph{calibration} (from projecting graph $\to$ physical
spectrum), \emph{embedding} (from embedding the graph into a
continuous space for measurement), and \emph{flow} (from
time-dependent dynamics).

The Bargmann graph supplies a worked example that refines this
classification.  The Coulomb $S^3$ graph has both calibration and
embedding $\pi$: calibration $\pi$ in Paper~2's $K = \pi(B + F -
\Delta)$~\cite{paper2_alpha}, embedding $\pi$ in the
electron-electron $1/r_{12}$ Slater integrals and the radial
normalization of hydrogenic functions.  The harmonic oscillator
graph has only embedding $\pi$ (the $\pi^{-3}$ Gaussian
normalization), and that embedding $\pi$ is never used by the
graph itself.

The refinement is therefore: \emph{calibration} exchange
constants are tied to nonlinear projections from second-order
Riemannian operators.  They are not a generic feature of quantum
discretization, but a specific feature of the Coulomb / $S^3$
construction (and, we expect, of any other system whose dynamical
symmetry group naturally acts on a Riemannian manifold rather
than a complex one).

%================================================================
\section{The harmonic oscillator rigidity theorem}
\label{sec:rigidity}
%================================================================

The Fock projection rigidity theorem of Paper~23~\cite{paper23}
states that the $S^3$ Fock projection $p_0 = \sqrt{-2 E_n}$ maps
a one-electron central-field Hamiltonian onto the free Laplacian
on $S^3$ if and only if the spectrum is $l$-degenerate within each
$n$-shell, which is unique to the Coulomb potential by
$\mathrm{SO}(4)$ symmetry.

The harmonic oscillator analog has the same structural form:

\begin{theorem}[Bargmann--Segal rigidity for the 3D HO]
\label{thm:rigidity}
The three-dimensional isotropic harmonic oscillator $V(r) =
\frac{1}{2} m \omega^2 r^2$ is the unique central-field
potential whose bound-state spectrum arises from the Euler
operator on the Hardy space $H^2(S^5)$ restricted to the
symmetric irreducible representations $(N, 0)$ of $\mathrm{SU}(3)$.
\end{theorem}

\begin{proof}[Proof.]
The argument has three independent classical inputs.

\emph{Step 1 (Jauch--Hill, 1940).}  Jauch and
Hill~\cite{jauchhill1940} proved that the three-dimensional
isotropic harmonic oscillator is the unique central-field
potential whose dynamical symmetry group includes
$\mathrm{SU}(3)$.  The proof constructs the explicit
$\mathfrak{su}(3)$ generators as bilinears in the position and
momentum operators (the angular momentum $\mathbf{L}$ and the
quadrupole tensor $Q_{ij} = \frac{1}{2}(p_i p_j / m\omega +
m\omega x_i x_j)$); the requirement that these close into
$\mathfrak{su}(3)$ under commutation forces $V(r) \propto r^2$.
Both directions of the equivalence
\[
  V(r) = \tfrac{1}{2} m\omega^2 r^2
  \;\;\Longleftrightarrow\;\;
  \mathrm{SU}(3)\text{ dynamical symmetry}
\]
are established by this construction; see also
Wybourne~\cite{wybourne1974}, Chapter~16, and
Iachello~\cite{iachello2006}, Chapter~4, for modern treatments.

\emph{Step 2 (Borel--Weil).}  The Borel--Weil
theorem~\cite{borelweil} identifies the symmetric irreducible
representations $(N, 0)$ of $\mathrm{SU}(3)$ with the spaces of
holomorphic sections of the line bundle $\mathcal{O}(N) \to
\mathbb{CP}^2$.  By restriction along the fibration $S^1 \to S^5
\to \mathbb{CP}^2$ (Eq.~\eqref{eq:s5-fibration}), this is the
same data as the homogeneous degree-$N$ subspace
$\mathcal{H}_N$ of the Hardy space $H^2(S^5)$
(Eq.~\eqref{eq:hardy-decomp}).  Hence the $(N, 0)$ irreducible
decomposition of an $\mathrm{SU}(3)$-equivariant Hilbert space
coincides exactly with the Hardy-space decomposition by
holomorphic degree.

\emph{Step 3 (Bargmann--Segal--Hall).}  The Bargmann
transform~\cite{bargmann1961}, made rigorous as a unitary
intertwiner of group representations by
Hall~\cite{hall1994}, establishes a unique (up to phase)
unitary equivalence between $L^2(\mathbb{R}^3)$ with the
Schr\"odinger representation of the harmonic oscillator and
$F^2(\mathbb{C}^3)$ with the Bargmann representation, under
which the harmonic oscillator Hamiltonian maps to
$\hbar\omega(\hat{N} + 3/2)$ (Eq.~\eqref{eq:H-bargmann}).
The Euler operator $\hat{N}$ is the unique self-adjoint operator
on $F^2(\mathbb{C}^3)$ that is diagonal on the monomial basis
with eigenvalue $a + b + c$ on $z_1^a z_2^b z_3^c$.

\emph{Combination.}  Suppose $V(r)$ is a central-field potential
whose bound-state spectrum arises from the Euler operator on the
Hardy-space sectors $\mathcal{H}_N$, with each $\mathcal{H}_N$
carrying the $(N, 0)$ irreducible representation of
$\mathrm{SU}(3)$.  By Step~2, the Hilbert space of $V(r)$ then
decomposes under $\mathrm{SU}(3)$ into the symmetric irreducible
representations $(N, 0)$, one per $N$, and the Hamiltonian is
linear in $\hat{N}$ on each.  By Step~3, this is exactly the
image of the harmonic oscillator under the Bargmann transform.
By the uniqueness of the Bargmann intertwiner (up to phase) and
of the Schr\"odinger representation, the original Hamiltonian
must be unitarily equivalent to the harmonic oscillator
Hamiltonian.  By Step~1, the only central-field potential with
this property is $V(r) = \frac{1}{2} m \omega^2 r^2$, completing
the converse.  $\hfill\square$
\end{proof}

\paragraph{Remark on the converse step.}
The crucial input to the converse direction is the biconditional
form of the Jauch--Hill theorem (Step~1): the algebraic closure
of the angular momentum and quadrupole generators into
$\mathfrak{su}(3)$ \emph{forces} $V \propto r^2$.  This is the
HO analog of the corresponding biconditional Fock--Bargmann
result for the Coulomb problem (the Runge--Lenz closure forces
$V \propto -1/r$).  Both biconditional results are standard
classical group theory~\cite{wybourne1974, iachello2006}.

The harmonic oscillator and the Coulomb problem are therefore
both rigid: the natural compact manifold and the natural
operator together uniquely determine the physical potential
through the closure conditions of the dynamical symmetry
algebra.  These two rigidity results are independent (Coulomb
uses $S^3$ and Laplace--Beltrami; HO uses $S^5$ Hardy space
and the Euler operator) and do not derive from each other, but
they have parallel structures and parallel proofs.  Together
they exhaust the two potentials that classical group theory
identifies as having full $\mathrm{SO}(4)$ and $\mathrm{SU}(3)$
dynamical symmetries respectively.

\subsection{Two-fermion entanglement rigidity (corollary)}
\label{sec:entanglement-rigidity}

The HO rigidity theorem of
Section~\ref{sec:ho-rigidity-theorem} has a direct
quantum-information corollary established empirically in
Paper~27~\cite{paper27_entropy}: \emph{the spatial 1-RDM
von Neumann entropy of the closed-shell two-fermion ground
state on the Bargmann--Segal graph is identically zero for
any central two-body interaction $V(r_{12})$.}

The mechanism is the Moshinsky--Talmi (MT) lab-to-relative
transformation~\cite{moshinsky_smirnov_1996}: any central
$V(r_{12})$ on the 3D HO basis preserves the total HO quantum
number $N_{\mathrm{tot}} = N_{\mathrm{rel}} + N_{CM}$.  The
two-body Hamiltonian therefore commutes with the one-body
Hamiltonian $H_{HO}$ at the operator level (both are
diagonal in $N_{\mathrm{tot}}$), and the closed-shell ground
state of $H_{HO}+V$ lies entirely within the lowest
$N_{\mathrm{tot}}$ sector.  For two fermions in the
$M_L=0, M_S=0$ singlet at the lowest shell ($N_{\mathrm{tot}}=0$),
that sector is one-dimensional --- a single Slater determinant
$|(0s)^{2}\rangle$ --- so the spatial 1-RDM has occupation
$(2,0,\ldots,0)$ exactly and the von Neumann entropy is zero.
Numerical verification at $N_{\max}\in\{2,3\}$ with the
Minnesota NN singlet potential at $\hbar\omega = 10$~MeV gives
$S_{\mathrm{full}}=0$ and $\lVert[H_{HO},V]\rVert_F /
\lVert H_{HO}\rVert_F < 10^{-15}$ (Paper~27,
Sec.~VII.A;~\cite{paper27_entropy}).

This is the structural dual of the corresponding statement
for the Coulomb $S^3$ basis, where Paper~27 establishes a
universal but \emph{nonzero} two-variable scaling
$S_B = A(\tilde{w}_B/\delta_B)^{\gamma}$ with
$\gamma\to 2$ at large $Z, n_{\max}$.  The Coulomb result is
nontrivial because $V_{ee} = 1/r_{12}$ does not preserve any
HO-like total-quanta quantum number on $S^3$; it couples
across shells, generating the second-order
Rayleigh--Schr\"odinger entanglement amplitude that
Paper~27 measures.  The HO basis kills this amplitude at
the operator level by the MT conservation law, completing
the universal/Coulomb-specific partition stated in the
abstract: the HO graph is $\pi$-free in eigenvalues, in the
Bargmann measure, and now also in two-fermion entanglement.

\subsection{Spinor-Lichnerowicz corollary: the operator-order discriminant}
\label{sec:spinor-lichnerowicz}

The HO rigidity theorem of Section~\ref{sec:ho-rigidity-theorem}
and the Coulomb/HO asymmetry of Section~\ref{sec:coulomb-ho-asymmetry}
established that calibration~$\pi$ (even-zeta content) enters from
second-order Riemannian operators on scalar bundles: the HO via
the Bargmann--Segal Euler operator on $H^2(S^5)$, and the Coulomb
system via the Fock projection to $-\Delta_{\mathrm{LB}}$ on~$S^3$.
The first-order/second-order distinction was identified as the
structural origin of the $\pi$-free property of the Bargmann graph
(first-order, complex-analytic, linear projection) versus the
calibration~$\pi$ of the Fock Dirichlet series (second-order,
Riemannian, nonlinear projection).

The Dirac-on-$S^3$ Tier~3 sprint (Track~T9) extends this result to
spinor bundles.  The squared Dirac operator on unit~$S^3$ satisfies
the Lichnerowicz identity $D^2 = \nabla^*\nabla + R/4$
\cite{lichnerowicz1963}, where $R = 6$.  Its spectral zeta evaluates
at every integer~$s$ to a two-term polynomial in~$\pi^2$ with
rational coefficients (Paper~18, Eq.~(T9)), with
\emph{no odd-zeta content} at any~$s$.

The mechanism is algebraic: squaring maps the first-order Dirac
sum $\sum g_n \cdot (\text{odd})^{-s}$ (which produces
$\zeta_R(\text{odd})$ at odd~$s$, as in Tier~1 Track~D3) to the
second-order sum $\sum g_n \cdot (\text{odd})^{-2s}$, which is always
$\mathbb{Q} \cdot \pi^{2s}$ by the Bernoulli formula.  The
odd-zeta content of~$D$ is structurally lost upon squaring.

\begin{corollary}[Operator-order transcendental discriminant]
The operator-order distinction (first-order vs.\ second-order) is
bundle-independent on round~$S^3$.  All second-order operators---scalar
Laplace--Beltrami, spinor connection Laplacian $\nabla^*\nabla$,
squared Dirac~$D^2$---produce exclusively $\pi^{\mathrm{even}}$
spectral-zeta content.  All first-order operators (Dirac~$D$) can
produce odd-zeta content ($\zeta_R(3), \zeta_R(5), \ldots$) via
half-integer eigenvalue sums.  The bundle type (scalar vs.\ spinor)
modulates rational coefficients and the eigenvalue/charge lattice
(integer vs.\ half-integer) but does not change the transcendental
class.
\end{corollary}

This corollary completes the universal/Coulomb-specific partition
stated in the abstract: the Bargmann graph is $\pi$-free because
it encodes a first-order operator (the Euler operator); calibration
$\pi$ is second-order; and this pattern persists on spinor bundles.
Cross-reference: Paper~18 \S IV (taxonomy closure);
Tier~1 Track~D3 ($\zeta(3)$ from first-order Dirac); T9 debug memo
(\texttt{debug/dirac\_t9\_memo.md}).

%================================================================
\section{Implications}
\label{sec:imp}
%================================================================

\subsection{For the fine-structure constant conjecture}
\label{sec:alpha}

Paper~2~\cite{paper2_alpha} conjectures the fine-structure
constant from the cubic equation $\alpha^3 - K \alpha + 1 = 0$
with $K = \pi (B + F - \Delta)$, where $B = 42$, $F = \zeta(2) =
\pi^2/6$, $\Delta = 1/40$ are spectral invariants of the Hopf
fibration $S^1 \to S^3 \to S^2$.  The numerical value $K =
137.036064$ matches $1/\alpha$ to $8.8 \times 10^{-8}$, with a
combinatorial $p$-value of $5.2 \times 10^{-9}$.  The combination
rule $K = \pi(B + F - \Delta)$ is empirical and the conjecture
remains conjectural.

The Coulomb / HO asymmetry developed in Section~\ref{sec:asymmetry}
provides structural support for one element of the formula: the
prefactor $\pi$.  By Theorem~\ref{thm:asymmetry}, calibration
$\pi$ of this type is tied to second-order Riemannian operators
with nonlinear projections.  The harmonic oscillator's
complex / first-order construction has no such $\pi$.  Therefore
$\pi$ in $K = \pi(B + F - \Delta)$ is structurally Coulomb-specific:
it comes from the spectral zeta values of $S^3$, which exist
because the Coulomb projection $E_n = -Z^2/(2 n^2)$ is nonlinear,
and there is no analog in any other dynamical-symmetry-classified
central potential.

This does not derive Paper~2's combination rule.  It explains why
$\pi$ appears in $K$, but does not explain the additive
combination $B + F - \Delta$ or the values of $B$, $F$, $\Delta$
themselves.  Paper~2 remains conjectural and its conjectural
status is unchanged by this paper.  The contribution here is to
sharpen the framing: any future derivation of the combination
rule must rely on the $\mathrm{SO}(4)$ / Riemannian features of
$S^3$ that distinguish the Coulomb problem from every other
central potential, not on the broader $\mathrm{SO}(3)$ angular
momentum structure that $S^3$ shares with $S^5$ and other
manifolds.

\subsection{For nuclear quantum simulation}
\label{sec:nuclear}

Paper~23~\cite{paper23} constructs nuclear shell-model qubit
Hamiltonians on the harmonic oscillator single-particle basis,
producing the deuteron and helium-4 qubit Hamiltonians via the
Minnesota nucleon-nucleon potential and Moshinsky--Talmi
brackets.  Theorem~\ref{thm:pi-free} of the present paper
implies that the single-particle structure of these Hamiltonians
is exactly rational at the graph level: the diagonal HO
energies, the orbital labeling, and the dipole transition
selection rules are all built from rational arithmetic on
integers.  Any floating-point error in nuclear Pauli coefficients
comes from the nucleon--nucleon interaction matrix elements
(which involve Gaussian integrals via Moshinsky--Talmi brackets),
not from the single-particle basis or the graph topology.  This
is a precision guarantee for nuclear qubit Hamiltonians built
within the GeoVac framework.

\subsection{The universal / Coulomb-specific partition completed}
\label{sec:partition}

Combined with Paper~22~\cite{paper22} (potential-independent
angular sparsity) and Paper~23~\cite{paper23} (Fock projection
rigidity theorem), the present paper completes a three-way
partition of the GeoVac framework's claims:

\paragraph{Universal across spherical fermion systems.}
The angular-momentum selection rules from Wigner $3j$ symbols
and the Gaunt integral; the ERI density bound (1.44\% at
$l_{\max} = 3$, independent of $V(r)$); the block-diagonal
structure of two-body operators; and the $O(Q^{2.5})$ composed
Pauli scaling.  These are the Paper~22 results.

\paragraph{Coulomb-specific.}
The $S^3$ Fock conformal projection; the calibration exchange
constant $\pi$; the Hopf bundle structure underlying the
$\alpha$ conjecture; the kinetic constant $\kappa = -1/16$ from
the $S^3$ Laplacian.  These are the Paper~23 / Track NH results
and the present paper.

\paragraph{Harmonic-oscillator-specific (this paper).}
The $S^5$ Hardy-space discretization; the linear-affine
projection $E_N = \hbar\omega(N + 3/2)$; the bit-exactly
$\pi$-free graph at every $N_{\max}$; the embedding $\pi^{-3}$
in the continuous Bargmann measure that the discrete graph
never invokes; the Bargmann--Segal rigidity theorem
(Theorem~\ref{thm:rigidity}).

The two natural geometries ($S^3$ for Coulomb and $S^5$ Hardy
for HO) are not in competition; they are the discrete encodings
of two physically distinct potentials with two distinct
dynamical symmetry groups.  The asymmetry between them is
structural and informs which features of the GeoVac framework
extend to nuclear, atomic, and molecular systems beyond the
electronic Coulomb problem.

%================================================================
\section{Conclusion}
\label{sec:conclusion}
%================================================================

We constructed a discrete graph encoding of the three-dimensional
isotropic harmonic oscillator on the holomorphic sector of $S^5$,
proved that the graph and its projection to physical energies are
bit-exactly $\pi$-free in exact rational arithmetic
(Theorem~\ref{thm:pi-free}), and developed the structural origin
of this fact in terms of the difference between Riemannian and
complex natural geometries (Theorem~\ref{thm:asymmetry}).  We
stated and proved the harmonic oscillator analog of the Fock
projection rigidity theorem (Theorem~\ref{thm:rigidity}), based on
the Jauch--Hill biconditional uniqueness of $\mathrm{SU}(3)$
dynamical symmetry, the Borel--Weil identification of $(N, 0)$
irreducibles with the holomorphic Hardy-space sectors, and the
Bargmann--Segal--Hall unitary intertwiner.

Calibration $\pi$ is Coulomb-specific.  It is the cost of mapping
a quadratic Riemannian eigenvalue to an inverse-square physical
energy, and it has no analog for the harmonic oscillator's linear
complex-analytic projection.  This refines Paper~18's exchange
constant taxonomy and sharpens Paper~2's framing of the
fine-structure constant conjecture: any geometric derivation of
$\alpha$ must rely on the $\mathrm{SO}(4)$ / Riemannian features
of the Coulomb problem, not on the angular momentum structure
shared with other central potentials.

Future work includes extending the discretization to anisotropic
harmonic oscillators (which break $\mathrm{SU}(3)$ to a smaller
subgroup, with corresponding modifications of the manifold), and
identifying the natural compact manifolds for other
dynamical-symmetry-classified central potentials such as the
Morse oscillator ($\mathrm{SU}(1, 1)$), the Kepler problem in
$n$ dimensions ($\mathrm{SO}(n + 1)$), and the P\"oschl--Teller
hierarchy.  Each such system should be classifiable by whether
its natural operator is first-order (no calibration exchange
constant) or second-order (calibration exchange constant
required), providing a structural taxonomy of where transcendental
constants enter the discrete-to-physical map.

\begin{thebibliography}{99}

\bibitem{fock1935}
V.~Fock,
``Zur Theorie des Wasserstoffatoms,''
\textit{Z.\ Phys.}\ \textbf{98}, 145 (1935).

\bibitem{bargmann1961}
V.~Bargmann,
``On a Hilbert space of analytic functions and an associated
integral transform,''
\textit{Comm.\ Pure Appl.\ Math.}\ \textbf{14}, 187 (1961).

\bibitem{segal1963}
I.~E.\ Segal,
``Mathematical problems of relativistic physics,''
in \textit{Proc.\ Summer Sem.\ Boulder, Colorado, 1960}
(American Mathematical Society, Providence, 1963).

\bibitem{hall1994}
B.~C.\ Hall,
``The Segal--Bargmann `coherent state' transform for compact Lie
groups,''
\textit{J.\ Funct.\ Anal.}\ \textbf{122}, 103 (1994).

\bibitem{follandstein1974}
G.~B.\ Folland and E.~M.\ Stein,
``Estimates for the $\bar{\partial}_b$ complex and analysis on the
Heisenberg group,''
\textit{Comm.\ Pure Appl.\ Math.}\ \textbf{27}, 429 (1974).

\bibitem{jauchhill1940}
J.~M.\ Jauch and E.~L.\ Hill,
``On the problem of degeneracy in quantum mechanics,''
\textit{Phys.\ Rev.}\ \textbf{57}, 641 (1940).

\bibitem{borelweil}
A.~Borel and A.~Weil, ``Repr\'esentations lin\'eaires et espaces
homog\`enes k\"ahl\'eriens des groupes de Lie compacts,'' as
exposited in J.-P.\ Serre, \textit{Repr\'esentations lin\'eaires et
espaces homog\`enes K\"ahleriens des groupes de Lie compacts}
(S\'eminaire Bourbaki, expos\'e 100, 1954).  Modern textbook
treatment in W.~Fulton and J.~Harris,
\textit{Representation Theory: A First Course} (Springer, 1991),
Lecture~23.

\bibitem{wybourne1974}
B.~G.\ Wybourne,
\textit{Classical Groups for Physicists}
(Wiley-Interscience, New York, 1974), Chapter~16.

\bibitem{iachello2006}
F.~Iachello,
\textit{Lie Algebras and Applications}
(Springer, Berlin, 2006), Chapter~4.

\bibitem{rowe1985}
D.~J.\ Rowe,
``Microscopic theory of the nuclear collective model,''
\textit{Rep.\ Prog.\ Phys.}\ \textbf{48}, 1419 (1985).

\bibitem{edmonds}
A.~R.\ Edmonds,
\textit{Angular Momentum in Quantum Mechanics}
(Princeton University Press, 1957).

\bibitem{varshalovich}
D.~A.\ Varshalovich, A.~N.\ Moskalev, and V.~K.\ Khersonskii,
\textit{Quantum Theory of Angular Momentum}
(World Scientific, Singapore, 1988).

\bibitem{paper0}
J.~Loutey,
``Geometric Packing and the Universal Constant $K = -1/16$,''
GeoVac Paper~0 (2026).

\bibitem{paper2_alpha}
J.~Loutey,
``The Fine Structure Constant from Spectral Geometry of the Hopf
Fibration,''
GeoVac Paper~2 (2026). [Conjectural.]

\bibitem{paper7}
J.~Loutey,
``The Dimensionless Vacuum: $S^3$ Conformal Geometry and the
Schr\"odinger Equation,''
GeoVac Paper~7 (2026).

\bibitem{paper14}
J.~Loutey,
``Structurally Sparse Qubit Hamiltonians from Hyperspherical
Lattices,''
GeoVac Paper~14 (2026).

\bibitem{paper18}
J.~Loutey,
``Spectral-Geometric Exchange Constants and the $\pi$-Free Graph
Principle,''
GeoVac Paper~18 (2026).

\bibitem{paper22}
J.~Loutey,
``Potential-Independent Angular Sparsity for Qubit Hamiltonians in
Spherical Fermion Systems,''
GeoVac Paper~22 (2026).

\bibitem{paper23}
J.~Loutey,
``Nuclear Shell Model Hamiltonians on the Hyperspherical Lattice,''
GeoVac Paper~23 (2026).

\bibitem{paper27_entropy}
J.~Loutey,
``Entropy as a Projection Artifact: One-Body Operators are
Entanglement-Inert on Sparse Lattices,''
GeoVac Paper~27 (2026).

\bibitem{moshinsky_smirnov_1996}
M.~Moshinsky and Yu.~F.~Smirnov,
\emph{The Harmonic Oscillator in Modern Physics}
(Harwood Academic, 1996).

\bibitem{lichnerowicz1963}
A.~Lichnerowicz,
\textit{Spineurs harmoniques},
C.~R.\ Acad.\ Sci.\ Paris \textbf{257}, 7--9 (1963).

\end{thebibliography}

\end{document}
